% section: results
\section{Results}
\label{sec:results}

\TODO{populated from recorded results files; no number in this section is typed by hand}

\subsection{The five claims, stated separately}
\label{sec:results:claims}

These are routinely conflated in this literature. They are not the same statement and are not
equally strong.

\begin{enumerate}
  \item \textbf{Official training reproduced.} \emph{No.} No published training recipe for this task
        is reproduced here; the closest published system used hardware and a teacher model
        unavailable to this work.
  \item \textbf{Official checkpoint evaluated.} \TODO{yes/no --- depends on whether RefChartQA's
        released checkpoint proves downloadable in Phase 4.4}
  \item \textbf{Matched student baseline trained.} \emph{Yes.} A direct-answer LoRA control is
        trained on matched data and evaluated under identical conditions, which is what separates
        ordinary domain adaptation from the contribution of grounding, typed plans and execution.
  \item \textbf{Before/after improvement achieved.} \TODO{the mandatory result}
  \item \textbf{Published baseline exceeded.} \TODO{only if genuinely matched --- same dataset
        version, split, metric implementation, input conditions and training-data regime. Otherwise
        this is context, not a win.}
\end{enumerate}

\subsection{Where the before/after comparison starts}
\label{sec:results:before}

The mandatory claim is an improvement over this project's own starting point, so that starting
point is stated first, on the same checkpoint, the same frozen slice and the same decoding as
everything that follows.

% tab_variant_selection — populated by: cdt-report --what variant_selection
\begin{table}[t]
\centering
\caption{Instruct versus Thinking on a frozen 200-question ChartQA validation slice, balanced
between human and machine questions. Thinking is selected only if it is at least two relaxed-accuracy
points better, produces at least 90\% valid compact JSON under the planned prompt, and has no more
than twice the median latency. The choice is frozen before any test split is opened.}
\label{tab:variant_selection}
\begin{tabular}{lrrr}
\toprule
Variant & Relaxed accuracy & Valid JSON (\%) & Median latency (s) \\
\midrule
\TODO{Instruct} & & & \\
\TODO{Thinking} & & & \\
\midrule
Threshold to switch & $+2.0$ & $\geq 90$ & $\leq 2\times$ \\
\bottomrule
\end{tabular}
\end{table}


Two of these numbers set the problem the fine-tune has to solve. The model answers half the
questions correctly while emitting a schema-valid record only about a third of the time, and its
own expression tree reproduces its own answer in roughly seven cases in ten. A record that does
not parse, or whose tree disagrees with its answer, is counted as a failure throughout this
report; none of them are quietly dropped.

\subsection{Headline table}

% tab_headline — populated by: cdt-report --what headline
% Grounding is a PER-SUBSET triple: the official evaluator produces no aggregate,
% and the published reference is the human subset alone (DECISIONS.md 0002).
\begin{table*}[t]
\centering
\caption{Headline results. Every cell carries a 95\% bootstrap confidence interval. Grounding is
reported per question subset because the official evaluator produces no aggregate; the published
reference figure applies to the human subset only, where $n=500$ and differences below roughly four
points are not distinguishable from sampling noise.}
\label{tab:headline}
\begin{tabular}{lccccc}
\toprule
& \multicolumn{2}{c}{ChartQA relaxed accuracy} & \multicolumn{3}{c}{RefChartQA AP@0.5} \\
\cmidrule(lr){2-3}\cmidrule(lr){4-6}
System & Human & Machine & Human & Machine & PoT \\
\midrule
Untouched model (zero-shot)     & \TODO{} & \TODO{} & \TODO{} & \TODO{} & \TODO{} \\
Direct-answer LoRA (control)    & \TODO{} & \TODO{} & \TODO{} & \TODO{} & \TODO{} \\
\textbf{Grounded plan + executor} & \TODO{} & \TODO{} & \TODO{} & \TODO{} & \TODO{} \\
\midrule
Published reference (3B, 1 epoch) & --- & --- & 32.83 & 59.28 & 39.32 \\
\bottomrule
\end{tabular}
\end{table*}


Every cell carries a bootstrap confidence interval, and every comparison states whether it is
matched. Grounding is reported as a per-subset triple, because the official evaluator produces no
aggregate and the published reference figure is the human subset alone
(Section~\ref{sec:setup:data}).

\subsection{The run itself}
\label{sec:results:curves}

% fig_training_curves
\begin{figure}[t]
\centering
% \includegraphics[width=\linewidth]{figures/fig_training_curves.pdf}
\caption{\TODO{caption for training_curves}}
\label{fig:training_curves}
\end{figure}


\subsection{Grounding, stratified by target size}

% tab_stratified — populated by: cdt-report --what stratified
\begin{table}[t]
\centering
\caption{\TODO{caption for stratified}}
\label{tab:stratified}
\begin{tabular}{lccc}
\toprule
\TODO{columns} \\
\midrule
\TODO{rows -- generated, never typed by hand} \\
\bottomrule
\end{tabular}
\end{table}

% fig_ap_by_area
\begin{figure}[t]
\centering
% \includegraphics[width=\linewidth]{figures/fig_ap_by_area.pdf}
\caption{\TODO{caption for ap_by_area}}
\label{fig:ap_by_area}
\end{figure}


The aggregate figure is the least informative view of this result. Roughly half of all targets are
smaller than one visual token at the input resolution used, and those two populations behave as
different problems. Per-bucket item counts are given alongside the scores: a bucket containing forty
items carries almost no information and is not presented as though it does.

\subsection{Zero-shot grounding: an unpublished measurement}

No zero-shot RefChartQA grounding figure for this model exists in the literature. Producing it is a
deliverable in its own right, and the interpretation was fixed in advance of measuring so that it
could not be rationalised afterwards: a value well below the published reference confirms the
expected headroom; a value near it means headroom is thinner than hoped; and a value \emph{above} it
is reported as a finding rather than treated as a failure, with the headline becoming
trained-versus-matched-baseline --- which is the mandatory result in any case.

\TODO{state the measured value and which branch was taken}
